\newcommand{\CLDate}{18. April 2025}

\newcommand{\CLCompanyHeader}{
\fontsize{28pt}{32pt}\fontspec{Noto Sans}\textbf{ADVANTECH}
}

\newcommand{\CLContactMe}{
\begin{tabular}{l @{}}
  \CLDate \bigskip\\
  Bastian Schmidt \\
  Basler Straße 64\\
  81476 München\\
  (+49) 15112939370%\\
  % basti@schmidt-electronics.de
\end{tabular}
}

\newcommand{\CLContactCompany}{
\begin{tabular}{@{} l}
  Thomas Kastner\\ 
  Senior Director Embedded Engineering | EIoT Europe\\ 
  Advantech Europe BV\\ 
  Industriestr. 15\\
  82110 Germering

\end{tabular}
}

\newcommand{\CLPosition}{
{ \fontsize{12pt}{14pt}\fontspec{Noto Sans}\textbf{Application Assignment}}\\
}

\newcommand{\CLContent} {

  { \fontsize{8pt}{10pt}\textit{Scenario 1 Description}}\\
  The company ABC Corp. wants to build a HMI frontend for their industrial
  robots, consisting of \textit{RSB-3720} running \textit{Yocto} and
  \textit{AUO G070VW01 7“ LVDS} display.\\
  They are asking if the display is supported or how to enable/activate it? Can you help?

  { \fontsize{8pt}{10pt}\textit{Scenario 1 Answer}}\\
  According to the datasheet (compare [1]), all respective versions of the \textit{RSB-3720}
  (e.g. CD, WD, QC, and WQ) feature 1 Single Channel or 1 Dual Channel 24-bit LVDS.
  The \textit{AUO G070VW01 7“ LVDS} is a single channel LVDS panel (compare [2]) and should therefore be
  compatible with the \textit{RSB-3720} from a hardware perspective.

  Moreover, the \textit{RSB-3720} datasheet lists optional accessories, and the requested
  AUO LVDS display is included in that list. Accordingly, I assume that it is
  supported. To enable/activate the display, it is important to ensure that the cable pinout
  matches the connector on the \textit{RSB-3720} board.

  In order to enable the display, the Linux image must include the appropriate driver.
  The Advantech meta repository (see query in [3]) lists instances of LVDS displays
  in combination with AUO, so I assume a corresponding driver is available.
  If the driver is present in the image, the display should work automatically when connected.

  \bigskip


  { \fontsize{8pt}{10pt}\textit{Scenario 2 Description}}\\
  ABC Corp. have decided to use a different LVDS based display, with
  integrated USB based touch controller.\\
  Is this possible / what would be the impact or effort?

  { \fontsize{8pt}{10pt}\textit{Scenario 2 Answer}}\\
  If a compatible driver for the LVDS display is available, then using a different panel with
  an integrated USB-based touch controller should be possible.

  A query of the Advantech meta repository shows examples of enabling touchscreen support
  for certain boards. However, it is essential that a driver for the new touch controller
  exists and is supported on the \textit{RSB-3720}.

  One direct impact would be that one of the available USB ports will be occupied
  by the display’s touch controller, and thus not available for other peripherals.



  \newpage
  { \fontsize{8pt}{10pt}\textit{Scenario 3 Description}}\\
  The product has been moved to mass production and is shipping in volume.
  Unfortunately, due to a component shortage in the market, the
  \textit{Ethernet PHY Qualcomm AR8035} that has been used on \textit{RSB-3720}
  originally is not available anymore. Advantech have redesigned
  \textit{RSB-3720} to a new HW revision, using the
  \textit{Realtek RTL8211FS PHY}.\\
  What is the impact on Advantech and on customer side? Is it possible
  to support both HW versions with one BSP version only?

  { \fontsize{8pt}{10pt}\textit{Scenario 3 Answer}}\\
  If any customer Linux images do not support the new \textit{Realtek RTL8211FS PHY},
  they may experience issues—especially if they rely on Ethernet functionality, which is very likely.

  Supporting both hardware versions with a single BSP image is a sound approach.
  To achieve this, the image must be compiled with drivers for both PHYs, and no relevant
  drivers or kernel modules should be blacklisted.

  Additionally, a query of the Advantech meta repository suggests that the Realtek RTL PHY
  uses a common and widely supported firmware.

  In this case, customer impact should be minimal. However, if customers use
  custom-built images that lack the new PHY driver, they may encounter issues.
  Therefore, it is advisable to notify customers proactively about the hardware change.


  \bigskip
  \bigskip
  \bigskip

  [1] Advantech RSB-3720 datasheet\\\url{https://advdownload.advantech.com/productfile/PIS/RSB-3720/file/RSB-3720_DS(080122)20220801131816.pdf}

  [2] AUO G070VW01.012 specifications\\\url{https://www.auo-lcd.com/products/auo-lcd-screen/7inch/588.html}
  
  [3] ADVANTECH-Corp/meta-advantech2: auo query\\\url{https://github.com/search?q=repo\%3AADVANTECH-Corp\%2Fmeta-advantech2\%20auo\&type=code}
  
  [4] ADVANTECH-Corp/meta-advantech2: touchscreen query\\\url{https://github.com/search?q=repo\%3AADVANTECH-Corp\%2Fmeta-advantech2\%20touchscreen\&type=code}
  
  [5] ADVANTECH-Corp/meta-advantech2: RTL query\\\url{https://github.com/search?q=repo\%3AADVANTECH-Corp\%2Fmeta-advantech2+RTL\&type=code}

}
